\documentclass[11pt]{article}
\usepackage[margin=1.1in]{geometry}
\usepackage{amsmath, amssymb, amsthm}
\usepackage{booktabs}
\usepackage{graphicx}
\usepackage{hyperref}
\usepackage{microtype}

\newtheorem{proposition}{Proposition}
\newcommand{\E}{\mathbb{E}}
\newcommand{\dkl}{D_{\mathrm{KL}}}

\title{Entropy-Gated Kelly Selection under Market Disagreement\\
with the Intrinsic GROUND Ratio}
\author{Peter~J.~Mercurio,~PhD}
\date{June 2026}

\begin{document}
\maketitle

\begin{abstract}
In this fourth paper of our series on the topic of portfolio optimization, we
take the entropic portfolio selection framework developed in the first three
papers \cite{mercurio2020repo, mercurio2020binary, mercurio2020gepo} into
production application: the systematic selection of short-dated vertical
credit spreads on S\&P~100 underlyings. We introduce the intrinsic GROUND
ratio, a per-candidate form of the Growth Rate Over UNiform Divergence
(GROUND) ratio of our second paper, which scores each candidate trade by its
Kelly expected growth under the trader's estimated outcome measure $P$,
discounted exponentially in the Kullback--Leibler divergence of $P$ from a
reference measure $Q$. The form is fully general in $(P, Q)$; here $P$ is the
realized-volatility-implied outcome distribution of the spread and $Q$ is the
market's implied-volatility counterpart. We establish the exact equivalence
of intrinsic GROUND selection to Hansen--Sargent multiplier preferences, and
show that to second order in the disagreement the entropic discount charges
each candidate the Kelly growth that would be forgone if the market's
measure were correct. We validate a midpoint-based
execution model on a large sample of same-day option prints, apply the
intrinsic GROUND ratio to six years of S\&P~100 credit spread data with a
further out-of-time period, and demonstrate how divergence-gated selection
outperforms growth-rate selection alone, which loses money under identical
conditions.
\end{abstract}

\noindent\textbf{Keywords:} intrinsic GROUND ratio; relative entropy;
Kullback--Leibler divergence; Kelly criterion; option portfolios; credit
spreads; variance risk premium; portfolio optimization

\section{Introduction}

Systematic selling of equity option premium is conventionally justified by the
variance risk premium: implied volatility (IV) exceeds subsequently realized
volatility (RV) on average, so options are on average sold rich
\cite{carrwu2009, bollerslev2009, bakshikapadia2003}. The practical
difficulty is not the existence of this premium but its selection. On any
given day, hundreds of candidate spreads carry positive model expected value
(EV), most of that EV is noise in the inputs, and an optimizer that maximizes
a noisy score reliably selects the noise. This paper, the fourth in our series
on the topic of portfolio optimization, addresses the selection problem with
the two-factor entropic structure developed across the first three papers
\cite{mercurio2020repo, mercurio2020binary, mercurio2020gepo}: growth rate as
the reward, relative entropy as the risk. A growth factor prices each
candidate as a Kelly bet \cite{kelly1956} under an RV-implied probability
measure. An entropy factor measures, in nats, how far that measure sits from
the IV-implied measure of the same spread, and discounts the growth factor
exponentially in that distance. The product is the intrinsic GROUND ratio.
Candidates are gated by a fixed threshold and the top five per entry day are
taken.

This paper also documents a change of price basis relative to earlier
revisions of the system. Earlier revisions scored candidates on last-trade
prices clamped to the bid--ask band. Live operation revealed that intraday
last prints on the two legs of a spread are asynchronous; pairing a stale
short-leg print near its ask with a stale long-leg print near its bid
fabricates credits that were never simultaneously quoted, in the worst cases
implying payoff odds above 10 on spreads whose guaranteed-fill credit was
negative at the same instant. The replacement basis, the quoted midpoint, is
validated empirically in Section~\ref{sec:execution} and is used identically
in the backtest (vendor end-of-day quotes) and in live operation (intraday
snapshots), making the ranked score the same object in both environments.

The remainder of the paper is organized as follows. Section~\ref{sec:lit}
reviews the related literature. Section~\ref{sec:growthrate} provides
background on the Kelly criterion and extends it to the three-state outcome
space of a vertical credit spread. Section~\ref{sec:relent} provides
background on relative entropy as a portfolio risk measure and recalls the
GROUND ratio. Section~\ref{sec:intrinsic} introduces the intrinsic GROUND
ratio, its choice of measures, and its theoretical properties. Section~\ref{sec:execution} presents the
midpoint execution model. Section~\ref{sec:results} applies the method to six
years of S\&P~100 credit spread data, including a worked example and a
comparison to Kelly criterion selection. Sections~\ref{sec:discussion}
and~\ref{sec:conclusions} discuss and conclude.

All logarithms are natural; entropies and divergences are in nats.

\section{Literature Review}
\label{sec:lit}

Modern portfolio selection begins with the mean-variance portfolio
optimization (MVPO) program of Markowitz \cite{markowitz1952}, whose limitations under non-Gaussian returns motivated
information-theoretic risk measures built on Shannon entropy \cite{shannon1948}.
Philippatos and Wilson \cite{philippatos1972} first proposed entropy as a
portfolio risk measure; Zhou, Cai, and Tong \cite{zhou2013} survey the
substantial literature that followed. The direct ancestors of the present
system are the return-entropy portfolio optimization (REPO) of Mercurio, Wu,
and Xie \cite{mercurio2020repo}, which replaces the variance objective with
portfolio entropy computed combinatorially; its extension to discrete-payoff
assets as discrete entropic portfolio optimization (DEPO), built on relative
entropy and the Kelly criterion, in
\cite{mercurio2020binary}, where the Kullback--Leibler divergence of a wager's
outcome distribution from the uniform reference is shown to be a convex risk
measure in the sense of \cite{follmer2002}, with the uniform distribution
playing the role of the minimum-risk wager; and the generalized entropic
portfolio optimization (GEPO) of \cite{mercurio2020gepo}, which extends DEPO to
intervals of continuous returns and applies it directly
to a wealth of option strategies, including the vertical credit spreads traded
here. The second paper also introduces the risk-adjusted ratio for comparing
growth rates of discrete-return portfolios, the Growth Rate Over UNiform
Divergence (GROUND) ratio \cite{mercurio2020binary}, in deliberate analogy to
the reward-to-variability ratio of Sharpe \cite{sharpe1966}: where the Sharpe
ratio measures excess return per unit of standard deviation, the GROUND ratio
measures excess expected growth rate per unit of relative entropy with respect
to the uniform distribution, with a minimum-risk portfolio playing the role of
the risk-free asset. The GROUND ratio is then used in \cite{mercurio2020gepo}
to measure the risk-adjusted performance of GEPO option portfolios.

Read as a series, the four papers trace one idea through increasing contact
with reality. REPO \cite{mercurio2020repo} establishes entropy as the risk
objective in place of variance. The second paper \cite{mercurio2020binary}
moves to relative entropy and the Kelly criterion for discrete payoffs and
builds the GROUND ratio. GEPO \cite{mercurio2020gepo} generalizes the
framework to intervals of continuous returns and to structured option
positions. The present paper is the production stress test: the framework
meets estimation noise and market microstructure on six years of data and a
live trading system. The adversary changes accordingly. In the first three
papers the case for entropy is distributional (variance is the wrong risk
functional for discrete and asymmetric payoffs); in the fourth, the case is
informational: the inputs themselves are noisy, the optimizer amplifies that
noise, and relative entropy, measured between two live probability models of
the same instrument, is the quantity that separates edge from artifact.

The growth factor descends from Kelly's information-rate interpretation of
log-optimal betting \cite{kelly1956}, its application to securities markets by
Thorp \cite{thorp1971}, and the capital growth theory consolidated in MacLean,
Thorp, and Ziemba \cite{maclean2011}. The connection between log-optimal
growth and information theory is classical \cite{coverthomas2006}: the growth
rate foregone by betting under a wrong measure is exactly a Kullback--Leibler
divergence \cite{kullback1951}, which motivates a divergence-based discount on
any growth estimate whose input measure is uncertain.

The IV-implied outcome probabilities used here are Black--Scholes--Merton
exercise probabilities \cite{blackscholes1973, merton1973}; the idea that
option prices encode a state-price density is due to Breeden and Litzenberger
\cite{breeden1978}. The empirical relation between implied and realized
volatility, and the bias of IV as a forecast, is documented by Christensen and
Prabhala \cite{christensen1998}; high-frequency measurement of realized
volatility follows Andersen, Bollerslev, Diebold, and Labys
\cite{andersen2003}. The variance risk premium that the growth factor
monetizes is documented at the index and single-name level by Carr and Wu
\cite{carrwu2009}, Bollerslev, Tauchen, and Zhou \cite{bollerslev2009}, and
Bakshi and Kapadia \cite{bakshikapadia2003}.

Exponential discounting in a relative entropy is the signature of entropic
risk measures \cite{follmer2002} and of robust control formulations in which
an agent guards against model misspecification within a relative-entropy
neighborhood of a reference measure \cite{hansensargent2008,
maccheroni2006}. In incomplete-market pricing, the minimal entropy martingale
measure of Frittelli \cite{frittelli2000} selects the pricing measure closest
in relative entropy to the physical one; the intrinsic GROUND discount inverts
this logic, treating large divergence between the trader's measure and the
market's as evidence against the trade rather than as a quantity to be
minimized. Finally, the execution model of Section~\ref{sec:execution} is
consistent with the finding of Muravyev and Pearson \cite{muravyev2020} that
effective option trading costs are substantially below quoted half-spreads
because executions cluster well inside the quoted band.

\section{Maximum Exponential Growth Rate}
\label{sec:growthrate}

\subsection{The Kelly Criterion}

We start by providing some background on the growth-rate objective. Consider
a repeated wager that pays $+b$ per unit staked with probability $p$ and
$-1$ with probability $q = 1 - p$. A bettor who stakes a constant fraction
$w$ of wealth on each round compounds wealth at the expected exponential
growth rate given by (\ref{eq:kelly-growth}),
\begin{equation}
G(w) \;=\; p \ln(1 + wb) + q \ln(1 - w),
\label{eq:kelly-growth}
\end{equation}
and the Kelly criterion \cite{kelly1956} selects the fraction $w^\ast$ that
maximizes (\ref{eq:kelly-growth}). For the two-outcome wager the maximizer
has the familiar closed form in (\ref{eq:kelly-frac}),
\begin{equation}
w^\ast \;=\; \frac{pb - q}{b},
\label{eq:kelly-frac}
\end{equation}
the bettor's edge divided by the odds. Thorp \cite{thorp1971} carried the
criterion from gambling to securities markets, and Cover and Thomas
\cite{coverthomas2006} give the information-theoretic treatment: the
log-optimal growth rate is a mutual information, and the growth foregone by
betting under an incorrect distribution is a Kullback--Leibler divergence.
This last fact is the seed of the present paper.

\subsection{The Credit Spread Outcome Model}
\label{sec:instruments}

Here we will extend the wager of (\ref{eq:kelly-growth}) to the vertical
credit spread. A bull put credit spread sells a put at strike $K_s$ and buys
a put at a lower strike $K_\ell$; a bear call credit spread sells a call at
$K_s$ and buys a call at a higher strike. Let $c$ denote the net credit
received per share and $w_d$ the strike width, so the maximum loss per share
is $m = w_d - c$ and the payoff odds are $b = c / m$. At expiry the spread
resolves into one of three states: full win (the underlying expires on the
safe side of the short strike, and the seller keeps $c$), full loss (the
underlying expires beyond the long strike, and the seller pays $m$), or a
partial outcome (the underlying expires between the strikes). We write the
outcome distribution as the triple $(p, r_0, q)$, with $p$ the probability of
the full win, $r_0$ the partial-zone mass, and $q$ the probability of the
full loss.

Writing the per-unit-risk payoff vector as $\{+b,\ +\alpha b,\ -1\}$, the
partial-zone payoff coefficient follows from the linear payoff between the
strikes under a uniform location assumption, as in (\ref{eq:alpha}),
\begin{equation}
\alpha(b) \;=\;
\begin{cases}
\dfrac{b - 1}{2b}, & b < 1,\\[4pt]
0, & b \geq 1,
\end{cases}
\label{eq:alpha}
\end{equation}
so that for $b < 1$ the partial state pays $\alpha b = (b-1)/2$ per unit
risked, the mean of $+b$ at the short strike and $-1$ at the long strike. For
$b \geq 1$ the geometric mean payoff of the partial zone is positive, and
setting it to zero is a conservative simplification that understates the
value of partial outcomes on high-odds spreads.

\subsection{Kelly Expected Value under the Trader's Measure}
\label{sec:growth}

For the three-state wager, the growth rate of (\ref{eq:kelly-growth}) becomes
(\ref{eq:kelly3}),
\begin{equation}
\ell(w) \;=\; p \ln(1 + wb) + r_0 \ln(1 + w\alpha b) + q \ln(1 - w),
\label{eq:kelly3}
\end{equation}
whose first-order condition reduces to a quadratic in $w$ (linear when
$\alpha = 0$, in which case $w^\ast = (pb - q)/(b(p+q))$, the three-state
generalization of (\ref{eq:kelly-frac})); we denote the maximizer $w^\ast$. The Kelly expected value of the candidate is then given
by (\ref{eq:kellyev}),
\begin{equation}
E \;=\; e^{\ell(w^\ast)} - 1,
\label{eq:kellyev}
\end{equation}
the per-trade geometric expected return at optimal sizing, and we define the
growth signal as $g = \ln E$. Because $\ell$ is the expectation of a concave
function of the outcome, a variance penalty is built into $E$ by Jensen's
inequality: of two candidates with equal arithmetic EV, the higher-variance
one has the lower $E$.

The probabilities $(p, r_0, q)$ entering (\ref{eq:kelly3}) are the trader's
estimated outcome measure $P$. For each leg we compute the Black--Scholes
$N(d_2)$ exercise probability \cite{blackscholes1973} using, in place of the
quoted IV, the underlying's realized volatility measured over the trailing
ten trading days (annualized close-to-close standard deviation),
tenor-matched to the one-to-four-day holding period. We denote this measure
$P_{\mathrm{rv}}$. The economic motivation for this choice, and for not using
the market's own probabilities here, is given in
Section~\ref{sec:measures}.

\section{Minimum Relative Entropy}
\label{sec:relent}

\subsection{Shannon Entropy and the Kullback--Leibler Divergence}

For the purposes of this paper, the risk of a candidate trade is a relative
entropy. We start by providing some background. The Shannon entropy
\cite{shannon1948} of a discrete distribution $P$ with support $\mathcal{X}$
is given by (\ref{eq:shannon}),
\begin{equation}
H(P) \;=\; -\sum_{x \in \mathcal{X}} P(x) \ln P(x).
\label{eq:shannon}
\end{equation}
The Kullback--Leibler divergence \cite{kullback1951}, also known as the
relative entropy of $P$ with respect to $Q$, is given by (\ref{eq:dkl}),
\begin{equation}
\dkl(P \,\|\, Q) \;=\; \sum_{x \in \mathcal{X}} P(x)
\ln\!\left(\frac{P(x)}{Q(x)}\right).
\label{eq:dkl}
\end{equation}
Gibbs' inequality (see MacKay \cite{mackay2003}) shows that the relative
entropy is always non-negative, as in (\ref{eq:gibbs}),
\begin{equation}
\dkl(P \,\|\, Q) \;\geq\; 0,
\label{eq:gibbs}
\end{equation}
with equality if and only if $P = Q$ almost everywhere.

\subsection{Relative Entropy as a Portfolio Risk Measure}

In \cite{mercurio2020repo} we argued that entropy provides a natural
alternative to variance for problems in which the return distribution is
discrete or otherwise ill-suited to mean-variance treatment, and in
\cite{mercurio2020binary} we showed that the relative entropy of a wager's
outcome distribution with respect to the uniform reference is a convex risk
measure in the sense of F\"ollmer and Schied \cite{follmer2002}, with the
uniform distribution itself playing the role of the minimum-risk wager.
Conventional risk measures (variance, value-at-risk, conditional
value-at-risk) are defined on continuous return distributions and become
degenerate or ill-conditioned on a small finite outcome space; relative
entropy fills exactly that gap.

\subsection{The GROUND Ratio}
\label{sec:groundratio}

To compare the growth rates of portfolios with different risk profiles, in
\cite{mercurio2020binary} we introduced the risk-adjusted ratio called the
Growth Rate Over UNiform Divergence (GROUND) ratio, in analogy to the
reward-to-variability ratio of Sharpe \cite{sharpe1966}. Let $U_m$ be the
$m$-state discrete uniform distribution. For a chosen portfolio $R_a$ and the
minimum-risk portfolio $R_b$ existing in the $m$-state event space, the GROUND
ratio $G_m$ is given by (\ref{eq:ground-original}),
\begin{equation}
G_m \;=\;
\frac{\E\!\left(G_a(\mathbf{w}_a) - G_b(\mathbf{w}_b)\right)}
     {\dkl(R_a \,\|\, U_m) - \dkl(R_b \,\|\, U_m)}
\;=\;
\frac{\E\!\left(G_a(\mathbf{w}_a) - G_b(\mathbf{w}_b)\right)}
     {H(R_b) - H(R_a)},
\label{eq:ground-original}
\end{equation}
where $G_a(\mathbf{w}_a)$ is the growth rate of the chosen portfolio with
weighting $\mathbf{w}_a$, $G_b(\mathbf{w}_b)$ is the growth rate of the
minimum-risk portfolio with weighting $\mathbf{w}_b$, $\dkl(R_a \| U_m)$ is
the relative entropy of the chosen portfolio with respect to $U_m$, and
$H(\cdot)$ is the Shannon entropy of (\ref{eq:shannon}). The GROUND ratio was
used in \cite{mercurio2020gepo} to measure the risk-adjusted performance of
GEPO option portfolios against Kelly criterion strategies.

\section{The Intrinsic GROUND Ratio}
\label{sec:intrinsic}
\label{sec:ground}

\subsection{From Comparative to Intrinsic Form}

The GROUND ratio of (\ref{eq:ground-original}) is a comparative score: the
ranking of a chosen portfolio depends on a reference portfolio $R_b$. In a
streaming daily-selection setting, where a small number of positions is
rebuilt from scratch every day from a changing candidate pool, this
cross-candidate coupling is awkward. To that end, here we introduce the
\emph{intrinsic} GROUND ratio, a per-candidate score that requires no
reference member. For a candidate $i$ with Kelly expected value $E$ from
(\ref{eq:kellyev}) computed under the trader's measure $P$, a reference
measure $Q$, and discount strength $k > 0$, the intrinsic GROUND ratio is
given by (\ref{eq:intrinsic}),
\begin{equation}
\Gamma_i \;=\; E \cdot e^{-k\,\dkl(P \| Q)}
\;=\; \left(e^{\ell(w^\ast)} - 1\right) e^{-k\,\dkl(P \| Q)}
\;=\; \exp\!\left(g - k\,\dkl(P \| Q)\right),
\label{eq:intrinsic}
\end{equation}
where $g = \ln E$ is the growth signal and $\dkl$ is the relative entropy of
(\ref{eq:dkl}). The first factor is the variance-adjusted expected wealth
gain per trade at Kelly-optimal sizing; the second factor is an entropic
discount taking values in $(0, 1]$. Their product is dimensionless and reads
as a divergence-discounted per-trade geometric expected return. Selection is
a gate followed by a cap: all candidates with $\Gamma_i \geq \theta$ qualify,
and the five highest-$\Gamma_i$ qualifiers per entry day are taken. The
canonical values used in this paper are $k = 10$ and $\theta = 0.05$
(Section~\ref{sec:sweep}). Across our parameter sweeps, selection under the
intrinsic form was indistinguishable from selection under the comparative
form of (\ref{eq:ground-original}).

\subsection{Choice of Measures}
\label{sec:measures}

The form (\ref{eq:intrinsic}) is fully general in the pair $(P, Q)$. The
published GROUND ratio takes $Q$ to be the uniform distribution $U_m$, under
which the entropic term reads as a risk measure on informational
concentration: a near-certain bet is a sharp informational claim that may not
survive estimation noise, and the convex-risk-measure property of
\cite{mercurio2020binary} formalizes this reading. In the present application
we instead take $Q$ to be the market's IV-implied outcome measure
$Q_{\mathrm{iv}}$ of the same spread, computed identically to
$P_{\mathrm{rv}}$ but under each leg's quoted IV.

The choice of measure in each factor is economically forced. The IV-implied
triple is risk-neutral: it is the measure under which every asset earns the
riskless rate, so a Kelly expected value computed under $Q_{\mathrm{iv}}$
against market prices is approximately zero for every fairly quoted spread by
construction. Implied volatility is a price, not a forecast. Any meaningful
estimate of expected value must therefore be taken under an estimate of the
physical measure, and the RV-implied triple is exactly that: the win
probability if the underlying continues to move as it has actually been
moving. The growth factor is accordingly computed under $P_{\mathrm{rv}}$,
and it is permissive: it prices the variance premium and is positive for most
candidates on most days. It is also a systematic overstatement, because part
of the gap between implied and realized volatility is earned compensation for
jump and crash risk that backward-looking realized volatility cannot see; a
portion of any $P_{\mathrm{rv}}$ expected value is insurance premium rather
than mispricing. The entropy factor is the correction: it charges the
estimate for its distance from the market's measure, suppressing exactly
those candidates whose apparent edge rests on claiming the market is most
wrong. The growth factor estimates expected value under the trader's physical
measure; the entropy factor prices that estimate's distance from the market's
risk-neutral measure, converting a raw variance-risk-premium claim into a
robustness-priced one. Neither measure alone suffices, and
Section~\ref{sec:gate} quantifies the division of labor: the ungated
$P_{\mathrm{rv}}$ expected value loses money in precise agreement with this
reading. Note that the same measure $P_{\mathrm{rv}}$ appears in both
factors: the divergence in (\ref{eq:intrinsic}) is weighted by the very
beliefs the growth factor bets, so it is largest in the states where the
claimed edge is concentrated and where being wrong costs the most.

\subsection{Theoretical Properties}
\label{sec:theory}

Define $J_k = g - k\,\dkl(P \| Q)$, the multiplier-preferences functional
of Hansen and Sargent \cite{hansensargent2008}, a member of the
variational-preferences class of Maccheroni, Marinacci, and Rustichini
\cite{maccheroni2006}; the exponential-in-divergence form of
(\ref{eq:intrinsic}) is likewise the signature of entropic risk measures
\cite{follmer2002}. The following two propositions make the connection
exact and give the gate its economic reading.

\begin{proposition}[Robust-preferences equivalence and the geometry of the
gate]
\label{prop:equiv}
For every candidate with $E > 0$, (i) $\Gamma_i = \exp(J_k)$ exactly, and
for any threshold $\theta > 0$ the acceptance rule
$\Gamma_i \geq \theta$ is identical to the rule $J_k \geq \ln\theta$;
in particular, intrinsic GROUND selection and Hansen--Sargent
multiplier-preference selection produce the same decisions and the same
rankings. (ii) In the $(\dkl, g)$ plane the acceptance region is the closed
half-plane
\begin{equation}
g \;\geq\; \ln\theta + k\,\dkl,
\label{eq:halfplane}
\end{equation}
a straight line with slope $k$ and intercept $\ln\theta$.
\end{proposition}

\begin{proof}
(i) By definition $g = \ln E$, so
$\exp(J_k) = e^{g} e^{-k\dkl} = E\,e^{-k\dkl} = \Gamma_i$, and the
equivalence of the two rules follows from the strict monotonicity of the
exponential. (ii) Take logarithms of $\Gamma_i \geq \theta$ and rearrange.
\end{proof}

The boundary of (\ref{eq:halfplane}) is the level set of the Lagrangian of
divergence-budgeted growth selection (maximize $g$ subject to
$\dkl \leq \varepsilon$), so $k$ admits a direct economic reading: it is
the shadow price of model disagreement, in nats of log-growth charged per
nat of relative entropy. The threshold $\theta$ sets the intercept, the
minimum divergence-free growth a candidate must offer.

\begin{proposition}[Second-order equivalence to the reverse divergence]
\label{prop:secondorder}
Let $P$ and $Q$ be distributions on a common finite outcome space with
$Q(x) > 0$ for all $x$, and write $\delta = \max_x |P(x) - Q(x)|$. Then
\begin{equation}
\dkl(P \| Q) \;=\; \tfrac{1}{2}\chi^2(P, Q) + O(\delta^3)
\;=\; \dkl(Q \| P) + O(\delta^3),
\label{eq:chi2}
\end{equation}
where $\chi^2(P, Q) = \sum_x (P(x) - Q(x))^2 / Q(x)$. Consequently
$J_k = g - k\,\dkl(Q \| P) + O(\delta^3)$.
\end{proposition}

\begin{proof}
Write $P(x) = Q(x) + \varepsilon(x)$ with $\sum_x \varepsilon(x) = 0$.
Expanding $\ln(1 + \varepsilon/Q)$ to second order,
$\dkl(P\|Q) = \sum_x (Q + \varepsilon)\left(\tfrac{\varepsilon}{Q} -
\tfrac{\varepsilon^2}{2Q^2}\right) + O(\delta^3)
= \tfrac{1}{2}\sum_x \tfrac{\varepsilon^2}{Q} + O(\delta^3)$, using
$\sum \varepsilon = 0$. The identical expansion of $\dkl(Q\|P)$ yields
the same leading term, which is $\tfrac{1}{2}\chi^2(P, Q)$.
\end{proof}

Proposition~\ref{prop:secondorder} supplies the Kelly-theoretic reading of
the gate. In the idealized proportional-betting setting of Kelly
\cite{kelly1956} and Cover and Thomas \cite{coverthomas2006}, an agent who
bets beliefs $P$ in a market whose true measure is $Q$ forgoes exactly
$\dkl(Q \| P)$ of expected log-growth per period: the reverse divergence is
the price of being wrong. By (\ref{eq:chi2}), the forward divergence inside
$\Gamma_i$ equals that price to second order in the disagreement, which for
the candidates traded here is small (median $\dkl$ of selected picks is
$0.0147$). The intrinsic GROUND ratio therefore discounts each candidate's
own-measure growth by $k$ times the growth that would be forgone if the
market's measure, rather than the trader's, were correct: $\Gamma_i$ is a
Kelly expected value charged an insurance premium against the market being
right, and $k$ is the number of such worst-case growth losses charged.
Beyond second order, and outside the proportional-betting idealization, the
identification is heuristic; the forward direction additionally weights the
disagreement by the trader's own beliefs, concentrating the penalty in the
states the growth factor actually bets. Because the exponential is monotone,
we prefer the multiplicative form throughout for its direct reading as a
per-trade risk-adjusted return.

\section{Midpoint Execution Model}
\label{sec:execution}

The price basis is the central empirical question for any spread-selling
backtest: the strategy's apparent edge is a small number obtained by
differencing quoted prices, so a biased basis manufactures or destroys the
edge before any selection logic runs. Three candidate bases were evaluated
against data.

First, last prints are disqualified. Across May--June 2026 vendor files,
restricted to contracts with same-day volume and one to four days to expiry
($n = 870{,}723$), 19.2\% of last prints lie outside the same day's closing
bid--ask band. Intraday the problem is worse: the two legs of a spread print
at different times against different underlying prices, and the difference of
two asynchronous prints is not a price. In live operation this fabricated
credits exceeding 90\% of strike width (implied odds $b > 10$) on spreads
whose short-bid-minus-long-ask credit was negative at the same instant.

Second, the midpoint is the empirical center of execution. On the same print
sample, the signed location of real trades inside the quoted band,
$(\mathrm{print} - \mathrm{mid})/\mathrm{halfwidth}$, has median $-0.005$:
trades center on the midpoint. This holds in every quoted-width bucket,
including quotes wider than 100\% of the option's own price (per-bucket mean
signed location between $-0.06$ and $+0.04$). Quoted width is defensive
posture, not execution cost: market makers fill well inside wide quotes,
consistent with the effective-spread evidence of \cite{muravyev2020}. This
finding rejects width-proportional slippage models. It is corroborated by the
only available real executions (five fills on 2026-05-28), which achieved
$0.73$ to $0.91\times$ the combo mid (mean $0.82$) with no relationship to
quoted width; the two widest-quoted names filled closest to mid.

The adopted model follows. Selection scores every candidate on the raw combo
mid, $c = \mathrm{mid}_{\mathrm{short}} - \mathrm{mid}_{\mathrm{long}}$,
identically in backtest (end-of-day quotes) and live operation (snapshot
quotes). Realized fills are modeled as a flat fraction of mid, with $0.80$ as
the reporting case (rounding the measured $0.82$ toward conservatism) and the
full ladder $0.80/0.85/0.90/0.95$ reported as a stress test
(Table~\ref{tab:fillstress}). Because the fill fraction is flat, it does not
change selection, only realized profit; the calibration rests on five
executions from a single session and is the softest number in this paper.

Finally, the backtest basis and the live basis agree. Live selection occurs
at 15:01 Eastern; backtest quotes are end-of-day. On 4{,}305 contracts matched
between live 15:01 snapshots and vendor end-of-day files across four
overlapping sessions, the ratio of end-of-day mid to 15:01 mid has median
$1.000$ (per-day medians $0.983$ to $1.029$). The end-of-day basis is an
unbiased proxy for the live entry basis.

\section{Credit Spread Selection with the Intrinsic GROUND Ratio}
\label{sec:results}

\subsection{Data and Selection Procedure}
\label{sec:design}

We use vendor end-of-day option data from January 2020 through December 2025
on S\&P~100 underlyings, with January through June 2026 reserved for
out-of-time evaluation. On each entry day (Monday through Thursday), for each
underlying, the system considers one-strike-wide vertical credit spreads
expiring the coming Friday (one to four days to expiry), both bull puts and
bear calls. The short leg is taken in an absolute delta band of
$[0.35, 0.65]$; both legs must show open interest of at least 100; the payoff
odds must satisfy $b \geq 0.30$; and the per-share maximum loss is capped at
\$5. Entry dates are restricted to exchange trading sessions: the vendor
files republish the prior session's quotes on market holidays with the days
to expiry decremented, and without this filter such stale rows admit entries
on days the market was closed. No trend filter and no index-volatility gate
are applied (both were
tested and rejected; each cost roughly a quarter of profit under earlier
configurations). Settlement uses the underlying's expiry-day close;
partial-zone wins are haircut to 50\% of intrinsic value to reflect pin and
assignment risk, while partial losses are taken in full. Flat sizing of one
contract per pick is the reporting case, with \$10{,}000 starting capital.

\subsection{A Worked Example}
\label{sec:example}

To illustrate the computation, consider the bull put credit spread on MRK
selected by the production system on Wednesday 2025-12-17, two days before
its Friday expiry. The underlying closed at \$99.18. The short leg is the
\$99 put (quoted $0.50 / 1.37$, midpoint $0.935$) and the long leg is the
\$98 put (quoted $0.09 / 0.51$, midpoint $0.300$), giving a mid credit of
$c = 0.635$ on a \$1 strike width, a maximum loss of $m = 0.365$, payoff
odds $b = 1.740$, and partial coefficient $\alpha = 0$ from
(\ref{eq:alpha}). Note that the short leg's quoted band ($0.87$) is nearly
as wide as the midpoint itself ($0.935$); by the evidence of
Section~\ref{sec:execution}, such width is defensive posture rather than
execution cost, and the midpoint remains the correct pricing basis.

The trailing ten-day realized volatility of MRK was $21.4\%$, against quoted
implied volatilities of $35.3\%$ (short leg) and $25.9\%$ (long leg). The
RV-implied outcome measure from the $N(d_2)$ computation is
$P_{\mathrm{rv}} = (p, r_0, q) = (0.543,\ 0.230,\ 0.227)$ and the
IV-implied reference is $Q_{\mathrm{iv}} = (0.523,\ 0.208,\ 0.269)$: the
trader's measure assigns more probability to the full win and less to the
full loss than the market prices, the signature of the variance risk premium
on this name on this day. The divergence from (\ref{eq:dkl}) is
$\dkl(P_{\mathrm{rv}} \| Q_{\mathrm{iv}}) = 0.0049$. With $\alpha = 0$
the first-order condition of (\ref{eq:kelly3}) is linear and gives
$w^\ast = (pb - q)/(b(p+q)) = 0.535$ and $\ell(w^\ast) = 0.1829$, so the
Kelly expected value from (\ref{eq:kellyev}) is
$E = e^{0.1829} - 1 = +20.1\%$. The intrinsic GROUND ratio from
(\ref{eq:intrinsic}) at $k = 10$ is then
$\Gamma_i = 0.2006 \times e^{-10 \times 0.0049} = 0.191$, which clears the
threshold $\theta = 0.05$, and the spread qualified. MRK closed on the
Friday at \$101.09, above the short strike, and the spread settled as a
full win.

\subsection{Backtest Performance}
\label{sec:headline}

Table~\ref{tab:headline} reports the canonical configuration on 2020--2025.

\begin{table}[h]
\centering
\caption{Canonical configuration ($k=10$, $\theta=0.05$), 2020--2025, one
contract per pick, fills at $0.80\times$ mid. Starting capital \$10{,}000.}
\label{tab:headline}
\begin{tabular}{lr}
\toprule
Trades & 2{,}250 \\
Final equity & \$50{,}798 \\
Total return & $+408\%$ \\
Sharpe (daily, annualized) & $+2.36$ \\
Sharpe (weekly, annualized) & $+2.44$ \\
Maximum drawdown & $-5.8\%$ \\
Win rate & $57.0\%$ \\
Mean profit per trade & \$18.13 \\
Outcome mix (win/partial/loss) & 1{,}061 / 397 / 792 \\
Median payoff odds $b$ of picks & 2.03 \\
Median $\dkl$ of picks & 0.0147 \\
\bottomrule
\end{tabular}
\end{table}

\begin{table}[h]
\centering
\caption{Fill-quality stress at the canonical cell, 2020--2025, one contract
per pick. Selection is identical in every row; only the realized fill
fraction changes. Sharpe and drawdown are not meaningful below the
break-even fill and are omitted there.}
\label{tab:fillstress}
\begin{tabular}{lrrrr}
\toprule
Fill & Final & Sharpe (wk) & Max DD & Win \\
\midrule
$0.50\times$ mid & $-\$24{,}346$ & --- & --- & $53.4\%$ \\
$0.60\times$ mid & \$856 & --- & --- & $54.6\%$ \\
$0.65\times$ mid & \$13{,}398 & $+0.36$ & $-28.2\%$ & $55.3\%$ \\
$0.70\times$ mid & \$25{,}897 & $+1.25$ & $-14.8\%$ & $56.0\%$ \\
$0.75\times$ mid & \$38{,}364 & $+1.95$ & $-8.9\%$ & $56.4\%$ \\
$0.80\times$ mid & \$50{,}798 & $+2.44$ & $-5.8\%$ & $57.0\%$ \\
$0.85\times$ mid & \$63{,}191 & $+2.78$ & $-3.9\%$ & $57.8\%$ \\
$0.90\times$ mid & \$75{,}550 & $+3.00$ & $-2.5\%$ & $58.5\%$ \\
$0.95\times$ mid & \$87{,}862 & $+3.14$ & $-1.7\%$ & $59.2\%$ \\
\bottomrule
\end{tabular}
\end{table}

Each five percentage points of fill quality is worth roughly \$12{,}000 of
final equity, and extending the ladder downward locates the break-even fill
at approximately $0.63\times$ mid: the strategy is profitable at
$0.65\times$ mid, near zero at $0.60$, and ruinous at $0.50$. Three
observations anchor the reported $0.80$ case well above that floor. The five
real executions of Section~\ref{sec:execution} filled between $0.73$ and
$0.91\times$ mid, every one of them above break-even; the print-location
study shows market executions centering on the midpoint itself; and the
reporting case of $0.80$ sits below every observed execution. The edge does
not require good fills, only fills no worse than $0.65\times$ the quoted
midpoint, a level the executable market has not approached in any
measurement we have made. Out of time, the
configuration frozen on 2020--2025 and applied to January--June 2026 produced
183 trades, final equity \$12{,}726 ($+27.3\%$), weekly Sharpe $+3.22$,
maximum drawdown $-6.3\%$, and a win rate of $57.9\%$, at the $0.80\times$
mid fill.

\subsection{The Selection Surface}
\label{sec:sweep}

The full $(k, \theta)$ grid was evaluated from a cached scoring of all
73{,}270 candidates, so every cell shares identical inputs and fills.
Table~\ref{tab:ksweep} reports the best threshold per $k$.

\begin{table}[h]
\centering
\caption{Best threshold per $k$ (by final equity), 2020--2025, $0.80\times$
mid fills, one contract per pick.}
\label{tab:ksweep}
\begin{tabular}{rrrrrrr}
\toprule
$k$ & $\theta$ & Trades & Final & Sharpe (wk) & Max DD & Win \\
\midrule
4  & 0.075 & 2388 & \$47{,}591 & $+2.34$ & $-6.4\%$ & 57.5\% \\
6  & 0.075 & 1917 & \$49{,}801 & $+2.52$ & $-5.3\%$ & 57.7\% \\
8  & 0.075 & 1730 & \$49{,}918 & $+2.54$ & $-4.1\%$ & 57.9\% \\
10 & 0.05  & 2250 & \$50{,}798 & $+2.44$ & $-5.8\%$ & 57.0\% \\
12 & 0.05  & 2097 & \$49{,}948 & $+2.44$ & $-5.2\%$ & 56.8\% \\
16 & 0.05  & 1863 & \$47{,}731 & $+2.40$ & $-3.4\%$ & 56.7\% \\
20 & 0.05  & 1730 & \$44{,}681 & $+2.27$ & $-3.8\%$ & 56.6\% \\
30 & 0.05  & 1493 & \$40{,}745 & $+2.18$ & $-5.2\%$ & 56.5\% \\
\bottomrule
\end{tabular}
\end{table}

The surface is a broad plateau: $k \in [6, 12]$ is within roughly 2\% of the
peak final equity, with the growth optimum at $(k, \theta) = (10, 0.05)$ and
the in-sample weekly Sharpe peak at $(8, 0.075)$. The canonical cell is the growth optimum itself,
$(k, \theta) = (10, 0.05)$: in a framework whose objective is expected
growth, the self-consistent selection criterion for the canonical
configuration is maximum realized growth, and we fix that criterion in
advance rather than selecting on any path statistic. Drawdown-efficient
corners exist on the grid (the $k = 16$ row trades roughly 6\% of final
equity for the shallowest drawdown shown), but maximum drawdown is a path
statistic, among the noisiest objects measurable on a single six-year path,
and we decline to select on it. Below $\theta = 0.03$ the strategy disintegrates (drawdowns of
$-74\%$ and worse; ungated, the full candidate stream barely breaks even),
while above $\theta = 0.09$ the 2026 out-of-time Sharpe turns negative.
Over-gating fails forward, a sign that the highest-$\Gamma_i$ tail is not where
the edge concentrates once the gate has done its work.

\subsection{Comparison to Kelly Criterion Selection}
\label{sec:gate}

As in the first three papers of this series, the benchmark is the Kelly
criterion. Selecting the top five candidates per day by the growth signal $g$
alone, with no divergence gate, is precisely Kelly criterion selection under
the trader's measure: it produces 5{,}745 trades and a total of
$-\$15{,}175$ ($-\$2.64$ per trade, win rate 55.0\%) under identical fills.
The identical pipeline with the $\Gamma_i$ gate produces $+\$40{,}798$
(Table~\ref{tab:headline}). The gate does not merely improve the growth
signal; it is the difference between a losing and a winning strategy. The
mechanism is the noise-selection effect described in the Introduction:
maximizing a noisy expected value selects the noise, and the cross-measure
divergence is an observable, per-candidate noise meter.

The same conclusion follows within the Kelly-selected candidates themselves.
Splitting the top-growth candidates at their median divergence
($\dkl = 0.0928$) separates a profitable half (total $+\$8{,}312$, mean
$+\$2.89$ per trade) from an unprofitable half (total $-\$23{,}487$, mean
$-\$8.18$), with Welch $t = 4.51$ ($p = 6.7 \times 10^{-6}$). High measured
disagreement between the trader's measure and the market's predicts realized
loss.

\section{Discussion}
\label{sec:discussion}

The results support a reading of the intrinsic GROUND ratio as a noise-aware
growth criterion rather than a return forecaster. The growth factor, taken
alone, is an aggressive estimate built from a realized-volatility measure
that is itself a forecastable but noisy object \cite{andersen2003,
christensen1998}; selecting its maxima harvests estimation error, and the
negative total of the ungated selection is the empirical face of that
selection effect. The divergence factor plays the role that relative entropy
plays in robust decision theory \cite{hansensargent2008} and entropic risk
measurement \cite{follmer2002}: it prices, per candidate, the distance
between the trader's measure and the market's, and the data show this
distance forecasts realized loss within the very candidates the growth factor
likes most. In the language of the GEPO framework \cite{mercurio2020gepo},
the production system replaces the portfolio-level entropy objective with a
per-candidate divergence gate, which is the form the idea takes when the
portfolio is rebuilt from scratch every day from a small fixed number of
positions. The change of reference measure (from the uniform distribution of
\cite{mercurio2020binary, mercurio2020gepo} to the market's IV-implied
measure) is the substantive theoretical step of this revision: it converts
the entropic term from a static concentration penalty into a dynamic,
instrument-specific disagreement measurement, while preserving the
robust-preferences reading $\ln\Gamma_i = J_k$ \cite{hansensargent2008,
maccheroni2006} unchanged.

The execution findings deserve emphasis because they generalize beyond this
strategy. Quoted width in short-dated single-name options is not an execution
cost schedule; executions center on the midpoint at every width, in line with
\cite{muravyev2020}, and last prints disagree with closing quotes a fifth of
the time even when the contract traded that same day. Any backtest of
multi-leg option strategies that prices legs from last trades, or that
penalizes fills proportionally to quoted width, is mismeasuring its own edge.

The softest quantities in the system are execution-level: the fill fraction
rests on five real executions, and selection on quoted mids retains residual
upward quote noise in its highest-ranked candidates. Both are measurable
going forward, and both bias in a known direction under the reported
$0.80\times$ mid fill, which is below every observed execution in the
calibration sample.

\section{Conclusions}
\label{sec:conclusions}

The intrinsic GROUND ratio operationalizes a simple division of labor: a
growth factor that prices the variance risk premium under the trader's
estimated physical measure, and an entropy factor that suppresses each
candidate exponentially in the measured disagreement between that measure and
the market's own. On a calibrated, basis-consistent execution model, the
gated strategy converts a losing Kelly criterion selection into $+408\%$ over
six years at a weekly Sharpe of $+2.44$ under conservative fills, with the
gate's contribution isolated and statistically strong. The configuration
carries forward the entropic portfolio optimization program of
\cite{mercurio2020repo, mercurio2020binary, mercurio2020gepo} into a
production setting in which the adversary is not distributional asymmetry in
returns but estimation noise in the inputs, and in which relative entropy,
measured between two live probability measures of the same instrument, is the
quantity that separates edge from artifact.

\section*{Materials and Methods}

Backtests use vendor end-of-day option chains (2020--2026) on S\&P~100
underlyings, a trailing ten-day realized-volatility table built from daily
closes, and per-leg quoted implied volatilities. The print-location study of
Section~\ref{sec:execution} uses all same-day-traded contracts with one to
four days to expiry in the May--June 2026 vendor files. Live operation runs
on snapshot quotes from a brokerage API at half-hour intervals, with the same
scoring code path as the backtest. All computations were performed in Python
with pandas and SciPy; the full parameter grid was evaluated from a single
cached scoring of all candidates so that every configuration shares identical
inputs and fills.

\section*{Abbreviations}

\begin{tabular}{ll}
GEPO & generalized entropic portfolio optimization \\
REPO & return-entropy portfolio optimization \\
DEPO & discrete entropic portfolio optimization \\
GROUND & growth rate over uniform divergence \\
MVPO & mean-variance portfolio optimization \\
EV & expected value \\
IV & implied volatility \\
RV & realized volatility \\
DTE & days to expiry \\
\end{tabular}

\begin{thebibliography}{99}

\bibitem{markowitz1952} H.~Markowitz, ``Portfolio Selection,''
\emph{Journal of Finance}, 7(1):77--91, 1952.

\bibitem{shannon1948} C.~E.~Shannon, ``A Mathematical Theory of Communication,''
\emph{Bell System Technical Journal}, 27:379--423, 623--656, 1948.

\bibitem{kullback1951} S.~Kullback and R.~A.~Leibler, ``On Information and
Sufficiency,'' \emph{Annals of Mathematical Statistics}, 22(1):79--86, 1951.

\bibitem{kelly1956} J.~L.~Kelly, Jr., ``A New Interpretation of Information
Rate,'' \emph{Bell System Technical Journal}, 35(4):917--926, 1956.

\bibitem{sharpe1966} W.~F.~Sharpe, ``Mutual Fund Performance,''
\emph{Journal of Business}, 39(1):119--138, 1966.

\bibitem{thorp1971} E.~O.~Thorp, ``Portfolio Choice and the Kelly Criterion,''
\emph{Proceedings of the Business and Economics Section of the American
Statistical Association}, 215--224, 1971.

\bibitem{philippatos1972} G.~C.~Philippatos and C.~J.~Wilson, ``Entropy, Market
Risk, and the Selection of Efficient Portfolios,'' \emph{Applied Economics},
4(3):209--220, 1972.

\bibitem{blackscholes1973} F.~Black and M.~Scholes, ``The Pricing of Options and
Corporate Liabilities,'' \emph{Journal of Political Economy}, 81(3):637--654,
1973.

\bibitem{merton1973} R.~C.~Merton, ``Theory of Rational Option Pricing,''
\emph{Bell Journal of Economics and Management Science}, 4(1):141--183, 1973.

\bibitem{breeden1978} D.~T.~Breeden and R.~H.~Litzenberger, ``Prices of
State-Contingent Claims Implicit in Option Prices,'' \emph{Journal of Business},
51(4):621--651, 1978.

\bibitem{christensen1998} B.~J.~Christensen and N.~R.~Prabhala, ``The Relation
between Implied and Realized Volatility,'' \emph{Journal of Financial
Economics}, 50(2):125--150, 1998.

\bibitem{frittelli2000} M.~Frittelli, ``The Minimal Entropy Martingale Measure
and the Valuation Problem in Incomplete Markets,'' \emph{Mathematical Finance},
10(1):39--52, 2000.

\bibitem{follmer2002} H.~F\"ollmer and A.~Schied, ``Convex Measures of Risk and
Trading Constraints,'' \emph{Finance and Stochastics}, 6(4):429--447, 2002.

\bibitem{andersen2003} T.~G.~Andersen, T.~Bollerslev, F.~X.~Diebold, and
P.~Labys, ``Modeling and Forecasting Realized Volatility,''
\emph{Econometrica}, 71(2):579--625, 2003.

\bibitem{bakshikapadia2003} G.~Bakshi and N.~Kapadia, ``Delta-Hedged Gains and
the Negative Market Volatility Risk Premium,'' \emph{Review of Financial
Studies}, 16(2):527--566, 2003.

\bibitem{mackay2003} D.~J.~C.~MacKay, \emph{Information Theory, Inference and
Learning Algorithms}, Cambridge University Press, 2003.

\bibitem{coverthomas2006} T.~M.~Cover and J.~A.~Thomas, \emph{Elements of
Information Theory}, 2nd ed., Wiley-Interscience, 2006.

\bibitem{maccheroni2006} F.~Maccheroni, M.~Marinacci, and A.~Rustichini,
``Ambiguity Aversion, Robustness, and the Variational Representation of
Preferences,'' \emph{Econometrica}, 74(6):1447--1498, 2006.

\bibitem{hansensargent2008} L.~P.~Hansen and T.~J.~Sargent, \emph{Robustness},
Princeton University Press, 2008.

\bibitem{carrwu2009} P.~Carr and L.~Wu, ``Variance Risk Premiums,''
\emph{Review of Financial Studies}, 22(3):1311--1341, 2009.

\bibitem{bollerslev2009} T.~Bollerslev, G.~Tauchen, and H.~Zhou, ``Expected
Stock Returns and Variance Risk Premia,'' \emph{Review of Financial Studies},
22(11):4463--4492, 2009.

\bibitem{maclean2011} L.~C.~MacLean, E.~O.~Thorp, and W.~T.~Ziemba (eds.),
\emph{The Kelly Capital Growth Investment Criterion: Theory and Practice},
World Scientific, 2011.

\bibitem{zhou2013} R.~Zhou, R.~Cai, and G.~Tong, ``Applications of Entropy in
Finance: A Review,'' \emph{Entropy}, 15(11):4909--4931, 2013.

\bibitem{mercurio2020repo} P.~J.~Mercurio, Y.~Wu, and H.~Xie, ``An
Entropy-Based Approach to Portfolio Optimization,'' \emph{Entropy}, 22(3):332,
2020.

\bibitem{mercurio2020binary} P.~J.~Mercurio, Y.~Wu, and H.~Xie, ``Portfolio
Optimization for Binary Options Based on Relative Entropy,'' \emph{Entropy},
22(7):752, 2020.

\bibitem{mercurio2020gepo} P.~J.~Mercurio, Y.~Wu, and H.~Xie, ``Option
Portfolio Selection with Generalized Entropic Portfolio Optimization,''
\emph{Entropy}, 22(8):805, 2020.

\bibitem{muravyev2020} D.~Muravyev and N.~D.~Pearson, ``Options Trading Costs
Are Lower than You Think,'' \emph{Review of Financial Studies},
33(11):4973--5014, 2020.

\end{thebibliography}

\end{document}
